% p.77 — الباب الرابع والثلاثين في قوانين التعليل وغيره
% Layout: \bab heading across the page, then twocol of numbered rules (circled numbers);
% each rule has the heading "قانون" on its own line followed by mostly unvocalised prose (only
% the example words carry ḥarakāt). rightcol: rules 1–6 up to «لان حرف العلة»; leftcol: the last
% line of rule 6 («ضعيف لا يحمل الحركات كما مر»), then rules 7–10 — split exactly where the
% lithograph breaks the column.
\mspage{77}
\bab{الباب الرابع والثلاثين في قوانين التعليل وغيره}
\begin{twocol}
\begin{rightcol}
\textbf{(1) قانون}\\
حروف العلة ثلاثة واو ياء الف ومعنى العلة سقم وضعف فلا يحمل الحركات الثلاثة الضَّمَّةُ والفتحة والكسرة الا نادرا
\par
\textbf{(2) قانون}\\
وباقي الحروف تسمى حروف الصحة ومعنى الصحة القوة فيحمل الحركات
\par
\textbf{(3) قانون}\\
الضَّمَّةُ نصف الواو مثل بُ بُوْ والكسرة نصف الياء مثل بِ بِيْ والفتحة نصف الالف مثل بَ بَا
\par
\textbf{(4) قانون}\\
الواو قائم مقام الضمتين والياء مقام الكسرتين والالف مقام الفتحتين
\par
\textbf{(5) قانون}\\
يناسب ان يكون ما قبل الواو مضموما وما قبل الياء مكسورا وما قبل الالف مفتوحا مثل عُوْدَ - عِيْدَ - عَادَ
\par
\textbf{(6) قانون}\\
ان كان حرف العلة متحركًا فيسكن - للتخفيف فيقال في بَيَعَ بَيْعَ وفي بُوِعَ بُوْعَ وفي بِيُعَ بِيْعَ لان حرف العلة
\end{rightcol}
\begin{leftcol}
ضعيف لا يحمل الحركات كما مر
\par
\textbf{(7) قانون}\\
ثم يبدل حرف العلة الفا ان كان ما قبله مفتوحا و وَاوًا ان كان ما قبله مضموما وياءً ان كان ما قبله مكسورا للمناسبة كما مر فيقال في بَيْعَ بَاعَ وفي بُيْعَ بُوْعَ وفي بِوْعَ بِيْعَ
\par
\textbf{(8) قانون}\\
ان كان حرف العلة مُتَحَرِّكًا فيسكن للتخفيف ثم يبدل مطابقا للحركة ما قبله للمناسبة فتقول في طَوُلَ طَوْلَ اولا ثم طَالَ ثانيا وفي بِوِعَ بِوْعَ ثم بِيْعَ وفي بُيِعَ بُيْعَ ثم بُوْعَ
\par
\textbf{(9) قانون}\\
اذا كان حرف العلة ساكنا اصليا و ما قبله مفتوحا فلا يقلب لِخِفَّةِ السُّكُوْنِ وَالْفَتْحَةِ مثل الْقَوْلِ وَالْبَيْعِ
\par
\textbf{(10) قانون}\\
اذا كان حرف العلة في آخر الاسم و ما قبله ساكنا فلا يحذف الحركة لمعرفة الاعراب مثل دَلْوٌ ظَبْيٌ - دَلْوًا ظَبْيًا - دَلْوٍ ظَبْيٍ
\par
\end{leftcol}
\end{twocol}
